From now on, let us f\/ix a singularity $x \in M$ of an inf\/initesimal automorphism $X \in \mathrm{inf}(\mathcal{M})$. We also choose a point $u \in \Gbdle$ in the f\/iber over $x$, and let $u_0 = \pi_+(u) \in \Gbdle_0$, likewise in the f\/iber over $x$. The remaining text is aimed at relating the local essentiality of $X$ near $x$, to invariant properties of the holonomy of $X$ at $x$, which is a one-parameter subgroup $h^t \subset P$:

\begin{Definition}[cf.~\cite{Frances3}, Section~6] \label{holonomy definition} Given $X$, $x$ and $u$ as above, \emph{the holonomy $h_u^t$ of $X$ at $x$ with respect to $u$} is def\/ined, for $t$ suf\/f\/iciently small, as follows: Let $\Phi_{\mathbf{X},t}(u)$, the integral curve of $\mathbf{X}$ through $u$, be def\/ined for $t \in (-\varepsilon,\varepsilon)$. Since $\mathbf{X}$ projects to $X$, $\mathbf{X}(u')$ is tangent to $\Gbdle_x$ for all $u' \in \Gbdle_x$, and hence all $\Phi_{\mathbf{X},t}(u)$ lie in $\Gbdle_x$. Then $h_u^t \in P$ is def\/ined by:
\[
\Phi_{\mathbf{X},t}(u) =: u.h^t_u.
\]
Since $h_u^{t+s} = h_u^th_u^s$ whenever both are def\/ined, $h_u^t = \mathrm{exp}(tX_{h,u})$ for some $X_{h,u} \in \p$, and we def\/ine~$h_u^t$ via this identity for all $t \in \R$.
\end{Definition}

Recall that, by def\/inition, $s_{\mathbf{X}}(u) = \omega(\mathbf{X}(u))$. Also, since $h_u^t = \mathrm{exp}(tX_{h,u})$, we have $X_{h,u} = (d/dt)_{\vert t=0} h_u^t$. By def\/inition, the integral curve $\Phi_{\mathbf{X},t}(u)$ satisf\/ies $\Phi'_{\mathbf{X}}(0) = \mathbf{X}(u)$. Hence, $s_{\mathbf{X}}(u) = \omega((d/dt)_{\vert t=0}(u.h_u^t)),$ and so by property (\ref{fund-vfs}) of the Cartan connection $\omega$, we have
\[
X_{h,u} = s_{\mathbf{X}}(u).
\]
In particular, this implies the following equivariance properties of $X_{h,u}$ and $h_u^t$ with respect to a change of the base point $u \in \Gbdle_x$, so it makes sense to speak of the holonomy $h^t$ of $X$ at $x$ as a conjugacy class of one-parameter groups in $P$:
\begin{gather*}
X_{h,u.p} = \mathrm{Ad}(p^{-1})(X_{h,u}), \qquad \forall \,\, p \in P;\\
h_{u.p}^t = p^{-1}h_u^tp, \qquad \forall \,\, p \in P.
\end{gather*}

The f\/irst part of relating essentiality of $X$ near $x$ to its holonomy, is the following:

\begin{Proposition} \label{necessary holonomy} Let $X \in \mathrm{inf}(\mathcal{M})$ have a singularity $x \in M$, and let $u \in \Gbdle_x$ be as above. If $X \in \mathrm{inf}(\sigma)$ for any locally defined Weyl structure $\sigma$, then $X$ has holonomy~$h_u^t$ conjugate under~$P$ to a one-parameter subgroup of~$G_0$ $($equivalently, $s_{\mathbf{X}}(u.p) = \mathrm{Ad}(p^{-1})(s_{\mathbf{X}}(u)) \in \g_0$ for some $p \in P)$.
\end{Proposition}

\begin{proof} Assume that $\sigma$ is any locally def\/ined Weyl structure in a neighborhood of the point $x$ with $X \in \mathrm{inf}(\sigma)$. From Lemma \ref{inessential automorphisms TFAE}, it follows that $\mathbf{X}_0 \in \mathrm{inf}(\Gbdle_0,\sigma^*\omega_{\leq})$, i.e.\ we have $\Lie_{\mathbf{X}_0} \sigma^*\omega_{\leq} = 0$. Then note that the identity (\ref{infinitesimal deformation identity}) from Lemma \ref{infinitesimal deformation} simplif\/ies, for $u \in \Gbdle_x$ and $u_0 := \pi_+(u) \in (\Gbdle_0)_x$, to give us the following two identities:
\begin{gather}
(\nabla^{\Abdle} s_{\mathbf{X}})(u) = 0; \label{Have0} \\
(\nabla^{\sigma} s_{\mathbf{X}_0})(u_0) = 0. \label{Need0}
\end{gather}
 To prove the claim in the proposition, we compute the identity (\ref{Have0}) in terms of $\sigma$, to show that~(\ref{Need0}) implies $s_{\mathbf{X}}(\sigma(u_0)) \in \g_0$. Since $\sigma(u_0)$, $u \in \Gbdle_x$, therefore $s_{\mathbf{X}}(\sigma(u_0)) = \mathrm{Ad}(p^{-1})(s_{\mathbf{X}}(u)) \in \g_0$ for some $p \in P$ as claimed.

 For the computation, note that in general, a Weyl structure $\sigma$ allows us to identify a section $s \in \Gamma(\Abdle)$ with $s \circ \sigma \in C^{\infty}(\Gbdle_0,\g)^{G_0} \isom \oplus_{i=-k}^k C^{\infty}(\Gbdle_0,\g_i)^{G_0}$. We will write $[s]^{\sigma} = (s^{\sigma}_{-k}, \ldots , s^{\sigma}_k) \in \oplus_{i=-k}^k C^{\infty}(\Gbdle_0,\g_i)^{G_0}$.

  Now consider any $Y \in T_xM$, and let $\mathbf{Y}_0$ be a local right-invariant vector f\/ield on $\Gbdle_0$ around $u_0 \in (\Gbdle_0)_x$, projecting onto $Y$ at $x$, and let $\mathbf{Y}$ be a local right-invariant vector f\/ield on $\Gbdle$ which extends the vector f\/ield $\sigma_* \mathbf{Y}_0$. Then by the chain rule, we have $\mathbf{Y}(s)(\sigma(u'_0)) = \mathbf{Y}_0([s]^{\sigma})(u'_0)$ for~$u'_0$ near~$u_0$ in~$\Gbdle_0$. Now compute from the def\/inition, for $s = s_{\mathbf{X}}$ as above:
\begin{gather*}
\big(\nabla^{\Abdle}_Y s\big)(\sigma(u_0))  = \mathbf{Y}(s)(\sigma(u_0)) + \{\omega \circ \mathbf{Y},s\}(\sigma(u_0))
  = \mathbf{Y}_0([s]^{\sigma})(u_0) + \{\sigma^*\omega \circ \mathbf{Y}_0,[s]^{\sigma}\}(u_0).
\end{gather*}

 Now we translate the last line into vector notation, where the top, middle and bottom components correspond, respectively, to the projection onto $\p_+$, $\g_0$ and $\g_-$ (denoted, as usual, with a~subscript). Note that from $\Pi(s) = X$, and since $X(x) = 0$, we have $s^{\sigma}_-(\sigma(u_0)) = 0$, and so we get the following reformulation of the left-hand side of~(\ref{Have0}):
\begin{gather}
\left(\begin{array}{c}
    \mathbf{Y}_0(s^{\sigma}_+)(u_0) \\
    \mathbf{Y}_0(s^{\sigma}_0)(u_0) \\
    \mathbf{Y}_0(s^{\sigma}_-)(u_0)\end{array}\right) +
\left(\begin{array}{c}
    \{\sigma^*\omega(\mathbf{Y}_0),[s]^{\sigma}\}_+(u_0) \\
    \{\omega_0(\mathbf{Y}_0),s_0^{\sigma}\}(u_0) + \{\omega_{-}(\mathbf{Y}_0),s^{\sigma}_+\}_0(u_0) \\
    \{\omega_{-}(\mathbf{Y}_0),s^{\sigma}_0\}(u_0) +
    \{\omega_{-}(\mathbf{Y}_0),s^{\sigma}_+\}_-(u_0)
     \end{array}\right). \label{Have0 wrt sigma}
\end{gather}

On the other hand, let us compute the identity (\ref{Need0}). The section $s_{\mathbf{X}_0} \in C^{\infty}(\Gbdle_0,\p^*)^{G_0}$ is given by
\[
s_{\mathbf{X}_0}(u'_0) = \sigma^* \omega_{\leq}(\mathbf{X}_0(u'_0)) = \omega_{\leq}(\sigma(u'_0))(\sigma_*(\mathbf{X}_0(u'_0))).
\]
Using the facts that $\sigma$ is a section of $\pi_+:\Gbdle \rightarrow \Gbdle_0$, and that $\mathbf{X}$ projects onto $\mathbf{X}_0$ via $\pi_+$, it follows that $\mathbf{X}(\sigma(u'_0)) - \sigma_*(\mathbf{X}_0(u'_0))$ lies in the kernel of $T_{\sigma(u'_0)} \pi_+$. In particular, this means we have:
\[
\omega_{\leq}(\sigma(u'_0))(\sigma_*(\mathbf{X}_0(u'_0))) = \omega_{\leq}(\mathbf{X}(\sigma(u'_0))),
\]
or equivalently, $s_{\mathbf{X}_0} = s^{\sigma}_- + s^{\sigma}_0$. Using this, a similar calculation to the one above gives:
\begin{gather}
(\nabla^{\sigma}_Y s_{\mathbf{X}_0})(u_0) =
    \left(\begin{array}{c}
    \mathbf{Y}_0(s^{\sigma}_0)(u_0) \\
    \mathbf{Y}_0(s^{\sigma}_-)(u_0)\end{array}\right) +
    \left(\begin{array}{c}
    \{\omega_0(\mathbf{Y}_0),s_0^{\sigma}\}(u_0) \\
    \{\omega_{-}(\mathbf{Y}_0),s^{\sigma}_0\}(u_0) \end{array}\right).
\label{Need0 wrt sigma}
\end{gather}

 Comparing the $\g_0$-components of (\ref{Have0 wrt sigma}) and (\ref{Need0 wrt sigma}), we see that if both terms vanish, we must have
\begin{gather*}
\{\omega_{-}(\mathbf{Y}_0),s^{\sigma}_+\}_0(u_0) := \mathrm{pr}_{\g_0}([\omega_{-}(\mathbf{Y}_0)(u_0),s^{\sigma}_+(u_0)]) = 0.
\end{gather*}
But we have $\g_- = \{\omega_{-}(\mathbf{Y}_0)(u_0) \, \vert \, Y \in T_xM\}$, and from this it follows that $s^{\sigma}_+(u_0) = 0$, by using the properties of $\vert k \vert$-graded semi-simple Lie algebras: We have the grading element $E \in \g_0$, which satisf\/ies $[E,Y_j] = jY_j$ for all $Y_j \in \g_j$. Now using the $Ad$-invariance of the Killing form $B$, we have, for any $Y \in \g_-$, and $0 < j \leq k$:
\[
B([Y,s^{\sigma}_j(u_0)],E) = B(Y,[s^{\sigma}_j(u_0),E]) = -jB(Y,s^{\sigma}_j(u_0)),
\]
and since $B$ induces an isomorphism $\g_j \cong (\g_{-j})^*$, this vanishes for all $Y \in \g_-$ only if $s^{\sigma}_j(u_0) = 0$ for all $j > 0$, i.e.\ only if $s(\sigma(u_0)) \in \g_0$, which is the claim of the proposition.
\end{proof}

Next we prove the converse to Proposition~\ref{necessary holonomy}:

\begin{Proposition} \label{sufficient holonomy} Let $X \in \mathrm{inf}(\mathcal{M})$ have a singularity $x \in M$, and let $u \in \Gbdle_x$, $u_0 = \pi_+(u) \in (\Gbdle_0)_x$. If the holonomy $h_u^t$ is conjugate under $P$ to a one-parameter subgroup of $G_0$ $($equivalently, $s_{\mathbf{X}}(u.p) = \mathrm{Ad}(p^{-1})(s_{\tilde{X}}(u)) \in \g_0$ for some $p \in P)$, then $X$ is an infinitesimal automorphism of some local Weyl structure~$\sigma$ around~$x$.
\end{Proposition}

\begin{proof} We will need the following ``exponential coordinates'' on $\Gbdle$ and $M$ around $u$ and $x$ respectively, which are induced by the Cartan connection $\omega$: For any $Y \in \g$, denote by $\hat{Y}$ the vector f\/ield on $\Gbdle$ which is determined by the identity, $\omega(\hat{Y}(u')) = Y$ for all $u' \in \Gbdle$. Def\/ining
\[
\mathcal{W}_u := \{ Y \in \g \, \vert \, \varphi_{\hat{Y},t}(u) \, \mathrm{is} \, \mathrm{def\/ined} \, \mathrm{for} \, 0 \leq t \leq 1 \},
\]
then there exist an open neighborhood $\mathcal{V}_u$ of $0 \in \g$ and an open neighborhood $V_u$ of $u \in \Gbdle$ such that the exponential map $\mathrm{exp}_u^{\omega}$ def\/ined on $\mathcal{W}_u$ is a dif\/feomorphism of $\mathcal{V}_u$ onto $V_u$, where by def\/inition:
\[
\mathrm{exp}_u^{\omega}: \  Y \mapsto \mathrm{exp}^{\omega}(u,Y) := \varphi_{\hat{Y},1}(u).
\]
Restricting $\mathcal{V}_u$ if necessary to a smaller neighborhood of zero, we get a dif\/feomorphism
\[
\overline{\mathrm{exp}}_u^{\omega} := \pi \circ \mathrm{exp}_u^{\omega}: \ \mathcal{V}_u^- \stackrel{\approx}{\rightarrow} U_x,
\]
where $U_x$ is a neighborhood of $x$ in $M$ and $\mathcal{V}_u^- := \mathcal{V}_u \cap \g_-$. Furthermore, for $V_u^- := \mathrm{exp}^{\omega}_u(\mathcal{V}_u^-)$, the restriction of the projection $\pi$ gives a dif\/feomorphism of $V_u^-$ onto $U_x$.

These exponential coordinates can obviously be used to def\/ine a local Weyl structure over~$U_x$, since they give a local section of $\pi: \Gbdle \rightarrow M$ on $U_x$. This gives us the local trivialisations $\pi^{-1}(U_x) \isom V_u^- \times P$ and $\pi_0^{-1}(U_x) \isom \pi_+(V_u^-) \times G_0$. Then we simply def\/ine $\sigma: \pi_0^{-1}(U_x) \rightarrow \pi^{-1}(U_x)$ by
\begin{gather}
\sigma: \ \pi_+(u').g_0 \mapsto u'.g_0; \label{exponential Weyl structure}
\end{gather}
for any $u' \in V_u^-$, $g_0 \in G_0$. This is by def\/inition a $G_0$-equivariant local section of $\pi_+: \Gbdle \rightarrow \Gbdle_0$, that is a Weyl structure. We will now show, assuming $h_u^t \subset G_0$, that the local f\/lows of $\mathbf{X} \in \VF(\Gbdle)$ commute with $\sigma$, i.e.\ we have $\Phi_{\mathbf{X},t} \circ \sigma = \sigma \circ \Phi_{\mathbf{X}_0,t}$ on $\pi_0^{-1}(U_x)$, for $t$ suf\/f\/iciently small so that both sides exist. By Def\/inition \ref{parabolic essential definition}, this implies $\mathbf{X} \in \mathrm{inf}(\sigma)$.

To show the commutativity of the local f\/lows of $\mathbf{X}$ with the Weyl structure $\sigma$ given by~(\ref{exponential Weyl structure}), we need the following general equivariance relation for an inf\/initesimal automorphism in exponential coordinates (cf.\ the proof of Proposition~4.2 of~\cite{Frances3}):
\begin{gather}
\Phi_{\mathbf{X},t}(\mathrm{exp}^{\omega}(u,Y)) = \mathrm{exp}^{\omega}(u,\mathrm{Ad}(h_u^t)(Y)).h_u^t. \label{holonomy inf aut equiv}
\end{gather}
The identity (\ref{holonomy inf aut equiv}) is based on the observation that we have $[\mathbf{X},\hat{Y}] = 0$ for any inf\/initesimal automorphism, and any $Y \in \g$ (this follows immediately from the def\/ining equation, $\Lie_{\mathbf{X}}\omega = 0$, of an inf\/initesimal automorphism). Hence, the f\/lows commute, $\Phi_{\mathbf{X},t} \circ \Phi_{\hat{Y},s} = \Phi_{\hat{Y},s} \circ \Phi_{\mathbf{X},t}$ whenever both sides are def\/ined, which together with equivariance of $\omega$ may be used to show that both sides of (\ref{holonomy inf aut equiv}) are given as the endpoint of the same integral curve through $\Phi_{\mathbf{X},t}(u) = u.h_u^t$.

Now consider an arbitrary point $\pi_+(u').g_0 \in \pi_0^{-1}(U_x)$, where $u' = \mathrm{exp}^{\omega}(u,Y) \in V_u^-$ for $Y \in \mathcal{V}_u^-$. Then $\sigma(\pi_+(u').g_0) := u'.g_0$, and we have, by $P$-equivariance of $\Phi_{\mathbf{X},t}$ and (\ref{holonomy inf aut equiv}):
\[
(\Phi_{\mathbf{X},t} \circ \sigma)(\pi_+(u').g_0) = \Phi_{\mathbf{X},t}(u'.g_0) = \mathrm{exp}^{\omega}(u,\mathrm{Ad}(h_u^t)(Y)).h_u^tg_0.
\]
But since $h_u^t \in G_0$, we have $\mathrm{Ad}(h_u^t)(Y) \in \g_-$ and for $t$ suf\/f\/iciently small we may also assume $\mathrm{Ad}(h_u^t)(Y) \in \mathcal{V}_u$ by continuity, so $\mathrm{exp}^{\omega}(u,\mathrm{Ad}(h_u^t)(Y)) \in V_u^-$. We also have $h_u^tg_0 \in G_0$, so $G_0$-equivariance of $\pi_+$ gives $\pi_+(\Phi_{\mathbf{X},t}(u'.g_0)) = \pi_+(\mathrm{exp}^{\omega}(u,\mathrm{Ad}(h_u^t)(Y))).h_u^tg_0,$ and hence combining the above gives:
\[
(\Phi_{\mathbf{X},t} \circ \sigma)(\pi_+(u').g_0) = \Phi_{\mathbf{X},t}(u'.g_0) = (\sigma \circ \pi_+ \circ \Phi_{\mathbf{X},t})(u',g_0).
\]
Finally, since $\Phi_{\mathbf{X}_0,t}$ is def\/ined on $\Gbdle_0$ via the relation $(\Phi_{\mathbf{X}_0,t} \circ \pi_+) = (\pi_+ \circ \Phi_{\mathbf{X},t})$, this shows that $(\Phi_{\mathbf{X},t} \circ \sigma) = (\sigma \circ \Phi_{\mathbf{X}_0,t})$ on~$\pi_0^{-1}(U_x)$.
\end{proof}

\subsection{Proof of Theorem~\ref{local essential dictionary}} \label{section4.3}

There are two parts to establishing Theorem~\ref{local essential dictionary}, and they are extensions of the arguments in the proofs of Propositions~\ref{necessary holonomy} and~\ref{sufficient holonomy}.

First, we establish the following claim, which is a strengthening of Proposition~\ref{necessary holonomy}: If $X$ is inessential in some neighborhood of $x$, then its holonomy $h_u^t$ is conjugate under $P$ to a one-parameter subgroup of $\mathrm{Ker}(\lambda) \subset G_0$ for a choice of scale representation $\lambda: G_0 \rightarrow \R_+$ (equivalently, $s_{\mathbf{X}}(u.p) = \mathrm{Ad}(p^{-1})(s_{\mathbf{X}}(u)) \in \mathrm{Ker}(\lambda') \subset \g_0$ for some $p \in P$).

The proof of this claim is analogous to the proof of Proposition \ref{necessary holonomy}. If $\sigma$ is exact, then we have the holonomy reduction $\overline{\Gbdle}_0 \subset \Gbdle_0$ to structure group $\mathrm{Ker}(\lambda) \subset G_0$, and denoting the resulting reduction by $\overline{\sigma}: \overline{\Gbdle}_0 \rightarrow \Gbdle$, the condition that $X$ is an \emph{exact} inf\/initesimal automorphism of $\sigma$ is (using Lemma \ref{inessential automorphisms TFAE}), in addition to the requirements computed in the proof of Proposition \ref{necessary holonomy}, that $\mathbf{X}_0(\overline{u}) \in T_{\overline{u}}\overline{\Gbdle}_0 \subset T_{\overline{u}}\Gbdle_0$ for all $\overline{u} \in \overline{\Gbdle}_0$ and that the resulting vector f\/ield $\overline{\mathbf{X}}_0$ on $\overline{\Gbdle}_0$ is an inf\/initesimal automorphism of the Cartan connection $\overline{\sigma}^*\omega_{\leq}$. This gives
\[
\big(\nabla^{\overline{\sigma}} s_{\overline{\mathbf{X}}_0}\big)(\overline{u}_0) = 0,
\]
where $\nabla^{\overline{\sigma}}$ denotes the tractor connection on the adjoint tractor bundle associated to $\overline{\Gbdle}_0$ by the adjoint representation on $\g_- \dsum \mathrm{Ker}(\lambda')$. Then by the same considerations, if we write $s^{\sigma}_0(\overline{u}_0) = s^{\overline{\sigma}}_0(\overline{u}_0) + z(\overline{u}_0)E_{\lambda}$, for $\overline{u}_0 \in (\overline{\Gbdle}_0)_x$, then this condition implies that
\[
0 = [\omega_-(\mathbf{Y}_0(\overline{u}_0)),z(\overline{u}_0)E_{\lambda}] = \sum_{j=1}^k jz(\overline{u}_0)\omega_{-j}(\mathbf{Y}_0(u_0)),
\]
for all $Y \!\in\! T_xM$, which can only happen if $z(\overline{u}_0) = 0$. Hence we must have $s(\sigma(\overline{u}_0)) \!\in\! \mathrm{Ker}(\lambda') \!\subset\! \g_0$.

Second, we make the following, additional claim under the setting of Proposition~\ref{sufficient holonomy}: If the holonomy $h_u^t$ of $X$ is conjugate under $P$ to a one-parameter subgroup of $\mathrm{Ker}(\lambda) \subset G_0$ for a choice of scale representation $\lambda: G_0 \rightarrow \R_+$ (equivalently, $s_{\mathbf{X}}(u.p) = \mathrm{Ad}(p^{-1})(s_{\mathbf{X}}(u)) \in \mathrm{Ker}(\lambda') \subset \g_0$ for some $p \in P$), then~$X$ is inessential on some neighborhood of~$x$.

This claim is also proved in the same way as the claim of Proposition \ref{sufficient holonomy}. The considerations of that proof also show that the local, equivariant section $\sigma$ of $\pi_+: \Gbdle \rightarrow \Gbdle_0$ over $U_x$ which was constructed using the exponential map, can be restricted to a map $\overline{\sigma}: (\overline{\Gbdle}_0)_{U_x} \rightarrow \Gbdle_{U_x}$ where the sub-bundle $(\overline{\Gbdle}_0)_{U_x} \subset (\Gbdle_0)_{U_x}$ is just given by $\pi_+(V_u^-) \times \mathrm{Ker}(\lambda)$ in the local trivialization, giving a locally exact Weyl structure. Finally, since $h_u^t$ is conjugate to a one-parameter subgroup of $\mathrm{Ker}(\lambda)$, it can be arranged that the restriction of $\mathbf{X}_0$ to $(\overline{\Gbdle}_0)_{U_x}$ is always tangent to this sub-bundle, so $X$ is locally inessential by Lemma \ref{inessential automorphisms TFAE}: Namely, changing $u \in \Gbdle_x$ if necessary, we may assume that $h_u^t \subset \mathrm{Ker}(\lambda)$ and (equivalently) $\omega(\mathbf{X}(u)) \in \mathrm{Ker}(\lambda')$. For any $u' \in V_u^- := \mathrm{exp}_u(\mathcal{V}_u^-) \subset \Gbdle_{U_x}$, the identity (\ref{holonomy inf aut equiv}) can be dif\/ferentiated to show that
\[
\mathbf{X}(u') = \frac{d}{dt}_{\vert t=0}\big(\mathrm{exp}^{\omega}(u,\mathrm{Ad}(h_u^t)(Y)).h_u^t\big),
\]
for some $Y \in \mathcal{V}_u^-$. Since $h_u^t \subset \mathrm{Ker}(\lambda)$, this shows that $\mathbf{X}(u')$ is the sum of a vector tangent to~$V_u^-$ with a vertical vector corresponding to an element of $\mathrm{Ker}(\lambda') \subset \g_0$. Thus, $\mathbf{X}_0(\pi_+(u')) \in T\overline{\Gbdle}_0$ and hence this holds for restriction of $\mathbf{X}_0$ to $(\overline{\Gbdle}_0)_{\vert U_x}$, since any point in the latter space is given by $\pi_+(u').k$ for some $u' \in V_u^-$ and some $k \in \mathrm{Ker}(\lambda)$.

